% ---- Title and author block ----
\title{\vspace{-1.5em}One Capability or Many?\\ Testing the Economic Validity of Frontier AI Evaluation}
\author{Louis Yiven Zhu\\
\small Department of Statistics, London School of Economics and Political Science}
\date{}
\maketitle

\begin{abstract}
\noindent
Frontier-model leaderboards increasingly report ``economic'' benchmarks that aim to measure performance
on professionally valuable work, such as software engineering, banking workflows, and other job-like
tasks. Whether these benchmarks capture a distinct economic capability, or merely re-express the general
capability signal that all benchmarks share as models improve, is a question of construct validity that
has not been tested directly. We address it on a pinned snapshot of 421 frontier model configurations,
of which 96 carry the complete twelve-benchmark battery that the structural analysis requires, framing
benchmark validity as a latent-variable measurement problem. We combine exploratory factor analysis
under an oblique rotation with a leave-one-benchmark-out predictive test, so that structural and
predictive evidence are assessed against pre-registered thresholds. A single factor accounts for 74.5\%
of common variance and is strongly associated with model release date ($R^2=0.51$); residualising on
release date reduces that share to 59.6\%, a 14.9-point drop whose bootstrap interval straddles our
15-point discriminance threshold, so the effect is indistinguishable from the threshold at this sample
size. Our dimensionality rule retains a single factor, so the registered distinctiveness test is
reported as not supported. When we extract more factors than that rule selects,
the economic benchmarks do separate, and the separation carries out-of-sample weight, improving prediction
of held-out economic scores over a single general-capability index on three of the four economic
benchmarks (pooled $\Delta\mathrm{MSE}=0.037$, 95\% bootstrap interval $[0.019, 0.055]$). Economic
benchmarks therefore carry limited, incremental validity. They refine a dominant progress factor at the
margin while remaining largely explained by it. All hypotheses are pre-registered, the failing and
near-threshold tests are reported as such, and every input is hash-pinned for exact reproduction.
\end{abstract}

\vspace{0.5em}
\noindent\textbf{Keywords:} construct validity; benchmark evaluation; large language models; frontier AI;
psychometrics; exploratory factor analysis; latent-variable models; economic capability; predictive
validity; pre-registration; reproducibility.
\vspace{0.5em}